Ask where a tree's mass comes from and intuition points downward: soil, water, nutrients drawn up through roots. This is almost entirely wrong. Trees are made of air—$\sim$95\% of their dry mass comes from atmospheric CO₂. Through photosynthesis, plants build themselves from carbon dioxide, converting invisible gas into solid wood, cellulose, and lignin using sunlight. Van Helmont's 1640s willow experiment demonstrated this: a tree gained 164 pounds ($\sim$74 kg) while the soil lost only two ounces ($\sim$60 g). Isotope labelling confirms the molecular accounting—carbon in wood comes from air, not earth or water. When trees burn, they simply return their borrowed carbon and sunlight to the atmosphere, completing a chemical cycle that temporarily crystallises air into living architecture.
